\section{Purchasing and Travel Oversight}
\textit{Reference $|$ Course Text: Chapters 15, 16}

As the PA, you are the approval authority for all team purchases and the facilitator for all travel planning. These responsibilities carry real financial consequences---unauthorized purchases create personal liability for students, and poorly planned travel wastes program funds. This section covers both domains.

% ═══════════════════════════════════════════════════════════════════
% PURCHASING
% ═══════════════════════════════════════════════════════════════════

\subsection{Purchasing: Your Role as Approval Authority}

Every purchase requires PA approval before the team spends money. This is not a rubber stamp---your approval confirms that the purchase is necessary, budgeted, and routed through the correct process. The CFO initiates the purchase; you authorize it.

\subsubsection{The Approval Workflow}

\begin{enumerate}[nosep,leftmargin=1.5em]
\item The team identifies a need and selects a component or material.
\item The CFO presents the purchase to you: what, why, how much, which vendor, and where it fits in the budget.
\item You verify the purchase is justified, within budget, and from an approved vendor.
\item The team submits the appropriate form to Capstone Purchasing (CP): Purchase Form, Shopping List, or Reimbursement Form.
\item CP processes the order and the team picks up from Labriola~204.
\end{enumerate}

\begin{bestpractice}
Establish a simple approval channel early---email or Slack message with the item, vendor, cost, and budget line. A quick ``Approved'' reply creates a paper trail. Do not let teams rely on verbal approvals alone; undocumented approvals lead to disputes.
\end{bestpractice}

\subsubsection{Reviewing the Budget Tracker}

The Budget Tracker is uploaded to Canvas monthly. Review it at least once per month and ask:

\begin{itemize}[nosep,leftmargin=1.5em]
\item Does actual spending match the approved budget categories?
\item Are there unexpected line items or categories trending over budget?
\item Has the team accounted for upcoming large purchases (custom PCBs, raw materials, shipping)?
\item Is the remaining budget sufficient to complete the project?
\end{itemize}

\begin{paaction}
Check the Budget Tracker at every sprint review during Sprints 8--11 (the build phase). This is when spending accelerates and overruns develop. A 30-second budget check during your meeting prevents surprises at FDR.
\end{paaction}

\subsubsection{Understanding the CP Workflow}

Capstone Purchasing (CP) processes all orders. Three forms exist on Canvas:

\begin{itemize}[nosep,leftmargin=1.5em]
\item \textbf{Purchase Form}: standard orders from approved vendors.
\item \textbf{Shopping List}: multiple items from the same vendor.
\item \textbf{Reimbursement Form}: student-made purchases (limited; requires your pre-approval, an itemized receipt, and no sales tax).
\end{itemize}

Teams must email CP (capstonepurchasing@mines.edu) with: team name/number, submission date, vendor name, and form used. CP processing takes 3--5 business days. Orders from unapproved vendors require a separate vendor review that may add 1--2 weeks.

\begin{tipbox}
Remind your teams that CP staff process hundreds of orders per semester. Complete, accurate forms get processed faster. If a team complains about slow processing, check whether they submitted the form correctly first.
\end{tipbox}

\subsubsection{Manufacturing Review Meeting Requirement}

Build Track teams must complete the \textbf{Manufacturing Review Meeting} before CP will process orders for build materials. This is a required checkoff on the purchasing form. If a team submits a materials order without this checkoff, CP will reject it.

\begin{paquestion}
\begin{itemize}[nosep,leftmargin=1.5em]
\item Has the Manufacturing Review Meeting been completed?
\item Does the BOM include all materials---including fasteners, consumables, and wiring?
\item Are long-lead items ordered or submitted to CP?
\end{itemize}
\end{paquestion}

\subsubsection{Common Purchasing Pitfalls}

Watch for these recurring problems across your teams:

\begin{redflag}
\textbf{Unauthorized purchases.} A student buys components with personal funds without PA approval, then submits a reimbursement request. This violates the purchasing workflow and may result in non-reimbursement. Reinforce at Sprint~7: no purchase without prior approval, no exceptions.
\end{redflag}

\begin{redflag}
\textbf{Missing tax exemption.} Mines is tax-exempt. Sales tax paid on purchases is \textbf{not reimbursable}. If a student buys locally without presenting the tax exemption certificate, the tax comes out of the team budget---or the student's pocket. Ensure the CFO has the certificate saved and accessible.
\end{redflag}

\begin{redflag}
\textbf{Late procurement.} Teams that delay ordering long-lead components (custom PCBs, specialty sensors, raw materials) past Sprint~8 risk not receiving parts before Spring Break. Items ordered after break have almost no schedule margin. Push teams to order critical components early and check procurement status at every sprint review.
\end{redflag}

\begin{warnbox}
If a team's spending is trending significantly over budget or you discover unauthorized purchases, escalate to the CF. Budget overruns and purchasing violations are course policy issues that may require intervention beyond your PA authority.
\end{warnbox}

\subsubsection{When to Escalate Budget Concerns}

Escalate to the CF when:

\begin{itemize}[nosep,leftmargin=1.5em]
\item A team's actual spending exceeds the approved budget by more than 15\%.
\item A student has made unauthorized purchases and is seeking reimbursement.
\item A team needs a budget increase that exceeds your comfort level.
\item A purchasing dispute arises between the team and CP that you cannot resolve.
\item A team is not submitting Budget Tracker updates on schedule.
\end{itemize}

% ═══════════════════════════════════════════════════════════════════
% TRAVEL
% ═══════════════════════════════════════════════════════════════════

\subsection{Travel: Facilitating Student Travel Planning}

Competition teams, field survey projects, and client site visits may require travel. As the PA, you facilitate the planning, approve the travel request, and serve as the safety backstop. The students handle logistics; you ensure the logistics are sound.

\subsubsection{The Go/No-Go Decision}

Before any travel spending begins, you must approve a \textbf{Go/No-Go decision} in coordination with the Capstone Leadership Team (CLT). The Go/No-Go considers:

\begin{itemize}[nosep,leftmargin=1.5em]
\item Is the trip essential to the project or competition?
\item Is the travel budget available and realistic?
\item Is the team prepared (prototype ready, competition paperwork complete)?
\item How many travelers are needed? Determine the \textbf{travel roster}---the minimum number of people required for the mission.
\end{itemize}

\begin{paaction}
Initiate the Go/No-Go conversation with CLT at least 8 weeks before the planned travel date. Late Go/No-Go decisions lead to expensive last-minute bookings and missed registration deadlines.
\end{paaction}

\subsubsection{Approving Travel Requests}

Once Go/No-Go is approved, the team's designated Travel Lead prepares a detailed budget. Before approving the travel budget, verify:

\begin{itemize}[nosep,leftmargin=1.5em]
\item Multiple lodging options with cost comparisons and links.
\item Airfare quotes per traveler (through FROSCH if required by Mines policy).
\item Ground transport plan: rental car, mileage for personal vehicles, or shared transportation.
\item Per diem estimates for all travelers for all travel days.
\item Registration and competition fees.
\item Equipment shipping costs (outbound and return).
\end{itemize}

\begin{bestpractice}
Require the Travel Lead to present the budget using the Budget Template on Canvas, with actual quotes---not guesses. Teams that present vague cost estimates invariably exceed their budget. If the budget is not detailed enough, send it back before approving.
\end{bestpractice}

\subsubsection{Understanding Mines Travel Policies}

Key policies you must enforce:

\begin{itemize}[nosep,leftmargin=1.5em]
\item \textbf{Carpooling required}: multiple personal vehicles for the same trip will not be reimbursed.
\item \textbf{Mileage}: reimbursed at the Mines rate, calculated from 1500 Illinois St., Golden CO 80401. The team must save the Google Maps route as a PDF.
\item \textbf{Fuel}: not reimbursable for personal vehicles (mileage rate covers fuel).
\item \textbf{Meals}: itemized receipts required. \textbf{No alcohol on any receipt.} A receipt containing alcohol is rejected \emph{in full}---no partial reimbursement.
\item \textbf{Post-trip receipts}: must be submitted to CP within one month of trip completion.
\item \textbf{Rental cars}: driver must be $\geq$21 years old with a valid license. Book through CP using the OneCard.
\end{itemize}

\begin{warnbox}
The alcohol receipt rule catches teams every semester. If any item on a receipt is an alcoholic beverage, the \textbf{entire receipt} is rejected. Remind teams before every trip: pay for alcohol separately on a personal card, or do not order alcohol on team travel.
\end{warnbox}

\subsubsection{Competition Team Travel Logistics}

Competition travel involves the largest budgets and most complex logistics in capstone. Key areas to oversee:

\textbf{Vehicle rental and ground transport.} Rental cars must be booked through CP using the OneCard. Verify the designated driver meets age and license requirements. For driving trips, ensure the team plans rest stops every 2--3 hours, limits driving to 10 hours per day, and switches drivers regularly.

\textbf{Hotel booking.} Two occupants per room to control costs. Book with a free cancellation policy. Verify proximity to the competition venue. Request late checkout on the last day if the event runs late.

\textbf{Registration and deadlines.} Competition registration often closes months before the event. Once Go/No-Go is approved, the Travel Lead must submit registration immediately. Track deadlines for tech inspection paperwork, safety forms, driver certifications, and other pre-event requirements.

\textbf{Equipment shipping.} If shipping the prototype or competition vehicle, the team must plan logistics well in advance: measure and weigh the shipment, get multiple carrier quotes, and ship at least 5 business days before the event. A prototype that arrives after the competition is a preventable disaster.

\begin{paquestion}
\begin{itemize}[nosep,leftmargin=1.5em]
\item Who is the designated Travel Lead, and do they have a detailed budget?
\item Have all travelers verified their IDs match their booking names?
\item What is the equipment shipping plan and timeline?
\item What is the contingency if a flight is cancelled or equipment is delayed?
\end{itemize}
\end{paquestion}

\subsubsection{Safety Considerations for Off-Campus Activities}

You are responsible for safety oversight of all off-campus activities, including travel and client site visits.

\begin{itemize}[nosep,leftmargin=1.5em]
\item \textbf{Itinerary sharing}: require the team to share the complete travel itinerary with you before departure. Every traveler should have your phone number.
\item \textbf{Emergency contacts}: for competition travel, identify the nearest hospital to the venue. Ensure every traveler has the PA's phone number, the campus emergency number, and the hotel address.
\item \textbf{Personal vehicle safety}: drivers must follow the Mines driving policy (available from EHS). Adequate sleep before long drives. Valid, current insurance.
\item \textbf{Client site visits}: before visiting a client site or field location, the team must assess site-specific hazards (construction, confined space, chemical, biological). Ask the client about required PPE and safety orientation before the visit.
\item \textbf{Weather and environment}: for outdoor or remote field work, verify the team has checked weather conditions, has cell coverage or a satellite communicator, carries a first aid kit, and has informed someone of their location and expected return time.
\end{itemize}

\begin{redflag}
A team plans an off-campus activity (site visit, field test, competition) without informing you. Any off-campus activity related to the project requires your awareness and, for travel, your explicit approval. If you discover a team visited a client site or conducted a field test without telling you, address it immediately.
\end{redflag}

\subsubsection{Required Documentation and Pre-Approval Process}

The travel planning timeline requires action well before the trip:

\begin{itemize}[nosep,leftmargin=1.5em]
\item \textbf{8--12 weeks before}: Go/No-Go conversation with CLT. Identify the travel need and designate the Travel Lead.
\item \textbf{6--8 weeks before}: Travel Lead develops the budget. Research lodging, airfare, and ground transport options. Register for the event.
\item \textbf{4--6 weeks before}: submit the travel budget to you and CP for approval. Book lodging and begin airfare booking. Confirm the travel roster.
\item \textbf{2--4 weeks before}: finalize all bookings. Prepare the equipment packing list. Verify all traveler documents (IDs, registration, competition paperwork).
\item \textbf{1 week before}: pre-trip meeting with the team. Distribute the itinerary. Verify all reservations. Confirm emergency contacts.
\item \textbf{Within 1 month after}: all receipts and reimbursement forms submitted to CP. Budget Tracker updated with actual travel costs.
\end{itemize}

For international travel, additional requirements apply: valid passports (6+ months remaining), registration with the Mines Office of Global Education (4+ weeks before departure), and verification of export control restrictions on any technical equipment.

\begin{paaction}
One week before departure, ask the Travel Lead to walk you through the pre-trip checklist: confirmed reservations (names match IDs), rental car and driver verification, competition registration and pre-event paperwork, equipment shipping status, emergency contact distribution, and printed copies of all itineraries and confirmations.
\end{paaction}


% ═══════════════════════════════════════════════════════════════════
